The \wikilink{s07--b001}{conjugation example} is the class equation.
The following definitions make its stabilizers explicit.

\begin{definition}\label{center-centralizer}
The \emph{center} of a group is
$Z(G)=\{z:zg=gz\text{ for all }g\in G\}$.
The \emph{centralizer} of $x\in G$ is
$C_G(x)=\{g:gx=xg\}$.
The \emph{conjugacy class} of $x$ is $\{gxg^{-1}:g\in G\}$.
\end{definition}
\begin{proposition}\label{center-centralizer-properties}
The center is a normal subgroup and each centralizer is a subgroup.
The conjugacy class of $x$ is an orbit of size $[G:C_G(x)]$.
\end{proposition}
\begin{proof}\label{center-centralizer-properties-proof}
Products and inverses of elements commuting with a fixed $x$ still
commute with $x$, giving the subgroup assertion. The center is the
intersection of these subgroups, and conjugation fixes every central
element. In the conjugation action the stabilizer of $x$ is exactly
$C_G(x)$; apply \wikilink{s06--b008}{orbit-stabilizer}.
\end{proof}
\begin{theorem}\label{class-equation}
For a finite group $G$, choose one representative from each noncentral
conjugacy class, forming a set $T$. Then
$|G|=|Z(G)|+\sum_{x\in T}[G:C_G(x)]$.
\end{theorem}
\begin{proof}\label{class-equation-proof}
Conjugacy classes partition $G$. The one-element classes are precisely
the central elements, and each remaining class has the size just
computed. Sum the sizes of the disjoint classes.
\end{proof}
\begin{proposition}\label{cyclic-central-quotient}
If $G/Z(G)$ is cyclic, then $G$ is abelian, without a finiteness assumption.
\end{proposition}
\begin{proof}\label{cyclic-central-quotient-proof}
Choose $a$ whose coset generates the quotient. Every element is
$a^mz$ for some integer $m$ and central $z$. Two such elements commute
because their powers of $a$ commute and their central factors commute
with every element.
\end{proof}
\begin{example}\label{abelian-conjugacy}
If $G$ is abelian, each centralizer is $G$ and all conjugacy classes
are singletons. The class equation then reduces to $|G|=|Z(G)|$.
\end{example}
